\section{Evaluation protocol}\label{sec:protocol}

\subsection{Benchmark, arms and what was run}
% source: EVAL:33-49, 58-64; RM:63-68; EXP:449-451, 502-504, 564, 670-676; D1 EXP:909.
The benchmark is derived from the sealed test split: 349 items over the seven tasks, in 78 groups, of which 271 are answerable (the reference commits) and 78 should-decline (the reference declines), in 61 and 17 groups respectively. Every item is in the JSON input format, which no training example uses. The protocol defines four arms (Table~\ref{tab:arms}); the pre-registered research question names the prompt-engineered orchestration baseline, so the comparisons in this paper are \code{arm1\_base\_orchestrated} against \code{arm2\_finetuned} (deviation D1: H1 had named the bare \code{arm0\_base}, of which only a 20-example probe exists, scored 0.5659, a biased head-of-file slice that is not comparable with a full run). If that probe were representative, the registered comparison would show a smaller gain than the orchestrated one does; the full bare-base run that would settle this has not been made. The orchestrated arm \code{arm3} was never run, so the interaction hypothesis H7 is only partly measured. Logos v0.0.2 was evaluated on \code{arm2} only: H9 does not ask whether fine-tuning helps, H8a had answered that for the architecture, and a base arm would have cost about 7 GPU hours (estimated) of Kaggle's 30-hour weekly quota. Fine-tuning's own effect on v0.0.2 is therefore not measured.

\begin{table}[t]
\centering\footnotesize
\caption{The evaluation arms and the runs in which each was evaluated. Source: \code{docs/evaluation.md} and the README.}
\label{tab:arms}
\begin{tabular}{L{44mm}L{28mm}L{32mm}L{44mm}}
\toprule
Arm & Model & Prompting & Evaluated in \\
\midrule
\code{arm0\_base} & base & none & no (a 20-example probe only) \\
\code{arm1\_base\_orchestrated} & base & \kleos{} orchestration & pilot, Hermes, Logos v0.0.1 \\
\code{arm2\_finetuned} & base $+$ LoRA adapter & controlled task context & all four runs \\
\code{arm3\_finetuned\_orchestrated} & base $+$ LoRA adapter & \kleos{} orchestration & no \\
\bottomrule
\end{tabular}
\end{table}

\subsection{Decoding and the reasoning split}
% source: configs/evaluation/kleos_policy_v006.yaml:40-52; kleos_logos_v002.yaml:8-23; EVAL:142-160; EXP:727-730; L2:136-145 (L-F7).
Every arm decodes greedily (temperature 0, top-$p$ 1, no sampling, repetition penalty 1) with batch size 1 and one seed, since greedy decoding makes the seed irrelevant. The generation budget is 512 new tokens, raised to 1,024 for Logos v0.0.2 because the worst case computed from the data, twice the longest training trace plus the longest benchmark answer, is $2 \times 298 + 199 = 795$ tokens. For the Reasoning model the completion is split on token id 35 before decoding and only the answer is graded; a completion whose trace never closes has an empty answer, is graded as such, and is marked as cut. The model card for that checkpoint recommends sampling at temperature 1; greedy decoding was kept so that every arm shares one decoding rule and so that the served model can be checked against the evaluation byte for byte (Section~\ref{sec:serving}). One consequence, finding L-F7, is that the stop reason reads ``stop'' whenever the trace closed even if the answer was then cut by the budget; 6 of the 349 v0.0.2 completions were cut that way, all after a closed trace, and were scored as generated.

\subsection{The grader}
% source: EVAL:108-124, 171-199, 250-268; src/kleos_models/evaluation/graders.py:230, 254; HER:559-569; L1:337; D6 EXP:909.
The composite grader \code{kleos\_policy} scores three components against the reference. The \emph{ranking} component is $0.6 \times \text{nDCG} + 0.4 \times \text{top-1 accuracy}$ against the reference ranking \citep{jarvelin2002ndcg}; nDCG was bounded after finding F1 so that a repeated item cannot score above a perfect ranking. The \emph{deciding factor} is read from the \code{deciding\_factor} key of a JSON answer, from a ``What decided it:'' line, or from the first allowed label the answer mentions, and compared with the reference label. \emph{Confidence} is read from an explicit boolean in a JSON answer or, otherwise, from a list of abstention phrasings taken from the \kleos{} corpus, and compared with whether the reference commits. The score is the mean of the components the reference supports. Whether the response is a JSON object (\code{format\_valid}) is reported as a sub-score and excluded from the score (deviation D6): under the format holdout a format failure must be distinguishable from a judgment failure. Kendall's $\tau$ \citep{kendall1938tau} and nDCG are reported as sub-scores. Three faithfulness heuristics (evidence coverage, citation precision and an unsupported-claim rate) are text-level heuristics, not entailment checks; the fabricated-citation heuristic flags 47.9\% of the gold test answers (finding H-F12) and evidence coverage is vacuous because no benchmark item names decisive evidence identifiers (H-F13), so they appear in this paper only beside those floors. Consistency across a group is measured by the agreement rate (the group reached one decision) and the correct-agreement rate (it agreed and was right). The configured grouping unit for the first runs was the scenario family, under which a perfect model scores only 0.600 because 6 of the 15 test families mix different correct labels (finding H-F11); from H8 on, consistency is also reported by the perturbation group, where a perfect model scores 1.000.

\subsection{Subsets and the unlearnable labels}
% source: D3 EXP:909 row; EXP:502-504, 668-672; LMD:323-324; SPEC:39; HER:363-367.
Deviation D3 was discovered during the analysis of the pilot: 78 of the 349 test items, the 78 should-decline items, expect a deciding-factor label that appears zero times in the v0.0.6 training targets (\code{request\_ambiguous} 30, \code{missing\_input} 25, \code{stale\_explicit\_conflict} 20, \code{insufficient\_separation} 3). The cause is a rendering rule: in v0.0.6, decline labels were printed only in JSON-format answers, and the format holdout put every JSON-format example in the test split, so 0 of the 124 training declines carry a label and 492 of the 820 training answers (60\%) have no ``What decided it:'' line at all. For a model trained on v0.0.6 the deciding-factor component is therefore capped at 0.7765 and the score at 0.9255 on the full benchmark, and the deciding factor is read on the answerable items as well as overall (0.816 on the 271 answerable items for the pilot, 0.6332 overall). This affects both arms of every v0.0.6 comparison identically, so H1 stays fair, but it is the reason the two cross-model hypotheses, H8b and H9, name the 271 answerable items as their primary population, with the 78 should-decline items as the first secondary row. v0.0.7 trains the four decline labels for the first time, and the should-decline row is the one it was built to move; H9 kept the answerable subset as primary on the stated ground that switching after the fact would ``pick the population most likely to flatter the result.''

\subsection{Statistics}
% source: EXP:818, 825-836; EVAL:215-217, 296, 318; metrics.py:330-366, 380, 395-436; EXP:489, 499-500, 665; XR:57-58; EXP:561-562, 765-766.
Differences are paired per item. A difference's 95\% interval comes from 2,000 bootstrap resamples with seed 42, taking the resampled differences at positions $\lfloor 0.025 B \rfloor$ and $\lfloor 0.975 B \rfloor$ of the sorted list; the two-sided $p$-value is twice the fraction of resamples on the far side of zero. When no resample crossed, the comparison tool prints $p = 0.0$ (H1) or $p < 0.0005$ (H8 onward, which is $1/B$); under the doubled one-sided fraction the two-sided bound is $2/B$, so this paper reports every such case as $p < 0.001$. A protocol amendment of 2026-09-24, adopted before any Logos training, changed the resampling unit: from H8 on, every interval resamples the 78 groups rather than the 349 items, because perturbations of one case are not independent and item-level intervals are too narrow (finding H-F14; \citealp{field2007clustered,miller2024errorbars}). Item-level intervals are still printed beside the cluster ones, no earlier result was re-decided, and intervals over fewer than 5 groups are reported as not estimable, which affects the 8-item \code{context\_prioritization} task. With 61 answerable and 17 should-decline groups, cluster inference is itself approximate; with few clusters it tends to over-reject \citep{cameron2008bootstrap}. Per-task rows are reported without correction for multiple comparisons and are not decisive. The decision rules are fixed in advance. For a fine-tuning hypothesis (H1, H8a) a task improves only if its interval excludes zero, and the hypothesis is supported if at least one task improves and none regresses. For a cross-model hypothesis (H8b, H9) the verdict on the paired cluster interval of the primary difference is \emph{better} if the interval lies entirely above zero, \emph{worse} if entirely below, \emph{equivalent} if entirely within $\pm 0.02$, and \emph{inconclusive} otherwise. Every result in this paper is reported as an estimate with its interval; a null is never read as equivalence \citep{altman1995absence}.

\subsection{Pre-registration}
% source: Table 4 from EXP:44-83, 136-146, 236-244, 273-281, 356-362, 484-500, 660-676; dates and commits from the git history (first commit fd87931, 2026-08-14) and the commit map XR:106-132; fixed protocol EXP:812-821.
Table~\ref{tab:hypotheses} lists the nine hypotheses with their metrics, decision rules, registration dates and outcomes. H1 to H7 and the fixed protocol (held-out splits only, greedy decoding identical across arms, a benchmark and graders fixed before the run, per-task reporting with no blended aggregate, recorded seeds, a leakage check before training) were committed on 2026-08-14 with the repository's first commit. H5 was re-scoped and H6 closed on 2026-09-15, after the pilot and before Hermes was trained. H8 and the clustered-interval amendment were committed on 2026-09-24, before any Logos training; deviation D11 was declared on 2026-09-29 before the arm it concerns ran; H9 was committed on 2026-10-06, before any Logos v0.0.2 training. In each case the commit that adds the entry is the registration and each run's manifest records the commit it ran at, so the order of registration and run is checkable from the public history (commit identifiers here are the post-rewrite ones; Appendix~\ref{app:artifacts} maps them to the identifiers the frozen records cite). Thirteen deviations, D1 to D13, are dated and logged with their effect on the result (Appendix~\ref{app:deviations}).

\begin{table}[t]
\centering\scriptsize
\caption{The pre-registered hypotheses. ``Registered'' is the date and post-rewrite commit of the entry; H5 to H7 carry no decision rule in the record. The record reports no interval for H3, although its rule names one. Source: \code{docs/experiments.md}, Status table and hypothesis entries.}
\label{tab:hypotheses}
\begin{tabular}{L{6mm}L{32mm}L{32mm}L{38mm}L{16mm}L{28mm}}
\toprule
Id & Question & Primary metric & Decision rule & Registered & Status \\
\midrule
H1 & Fine-tuning beats the prompt-engineered baseline & per-task grader score on the held-out test split (registered as arm0 vs arm2; run as arm1 vs arm2, D1) & improvement counts only if the paired bootstrap 95\% CI excludes zero; per task, no blended aggregate & 2026-08-14, \code{fd87931} & Supported on the pilot (2026-09-15, with D1 and D2); replicated by Hermes (2026-09-22) \\
H2 & Gains generalize out of distribution & \code{ood\_score} beside \code{in\_distribution\_score}; generalization gap & gains generalize if the OOD delta CI excludes zero & 2026-08-14 & Not measurable: the v0.0.6 benchmark is entirely out of distribution (2026-09-15) \\
H3 & Decisions survive irrelevant perturbations & \code{correct\_agreement\_rate}; flips per perturbation kind & consistency improves if the delta CI excludes zero & 2026-08-14 & Improved, still poor (pilot 2026-09-15; reproduced exactly by Hermes 2026-09-22); no interval computed \\
H4 & No loss of general capability & \code{relative\_delta} on a fixed capability suite & regression if $< -0.01$ with a CI excluding zero & 2026-08-14 & Not run \\
H5 & The effect depends on scale & --- & --- & 2026-08-14; re-scoped 2026-09-15 & Not run; re-scoped to two Mistral models \\
H6 & Transfer across model families & --- & --- & 2026-08-14; closed 2026-09-15 & Closed: not testable with one model family \\
H7 & Fine-tuning and orchestration interact & --- & --- & 2026-08-14 & Partially measured: arm1 and arm2 only, in three runs (2026-10-01) \\
H8a & Fine-tuning helps Logos v0.0.1 & per-task \code{kleos\_policy}, Logos arm1 vs arm2 & a task improves only if its cluster CI excludes zero; supported if at least one improves and none regresses & 2026-09-24, \code{81e3e9a} & Supported (2026-10-01) \\
H8b & Logos v0.0.1 is better than Hermes (primary) & mean \code{kleos\_policy} on the 271 answerable items, paired & better / worse / equivalent ($\pm 0.02$) / inconclusive, by the cluster CI of Logos $-$ Hermes & 2026-09-24, \code{81e3e9a} & Inconclusive (2026-09-29) \\
H9 & Logos v0.0.2, trained to think, is better than Hermes (primary) & as H8b & as H8b & 2026-10-06, \code{6d7e450} & Supported: better, $+0.0409$, cluster 95\% CI $+0.0040$ to $+0.0780$ (2026-10-06) \\
\bottomrule
\end{tabular}
\end{table}
